% =============================================================================
%  Full project report
%  Probabilistic Motor Subtype Progression Model for Parkinson's Disease
%
%  Figures are read from outputs/; regenerate with scripts/01 .. 07 then
%  compile twice with pdflatex.
% =============================================================================
\documentclass[11pt,a4paper]{article}

\usepackage[T1]{fontenc}
\usepackage[utf8]{inputenc}
\usepackage[margin=2.4cm]{geometry}
\usepackage{booktabs}
\usepackage{longtable}
\usepackage{array}
\usepackage{xcolor}
\usepackage{graphicx}
\usepackage{enumitem}
\usepackage{caption}
\usepackage{amsmath}
\usepackage{amssymb}
\usepackage{float}
\usepackage[hidelinks]{hyperref}

\graphicspath{
  {../outputs/transitions/figures/}
  {../outputs/model/figures/}
  {../outputs/hbayes/figures/}
}

\definecolor{warn}{HTML}{B7770B}

\newcommand{\code}[1]{\texttt{\small #1}}
\newcommand{\arr}{$\rightarrow$}

\setlist[itemize]{itemsep=2pt,topsep=4pt}
\setlist[enumerate]{itemsep=2pt,topsep=4pt}
\captionsetup{font=small,labelfont=bf}

\newenvironment{keypoint}
  {\begin{center}\begin{minipage}{0.94\textwidth}%
   \vspace{4pt}\hrule\vspace{7pt}\small}
  {\vspace{7pt}\hrule\vspace{4pt}\end{minipage}\end{center}}

\title{\textbf{A Probabilistic Model of Motor Subtype Progression\\
       in Parkinson's Disease}\\[0.7em]
       \large Project Report}
\author{Parv Patel \quad\textbullet\quad Mayank Tewatia\\[0.4em]
        \small Supervisor: Prof.\ K.~R.\ Sahasranand\\
        \small Department of Data Science, IIT Palakkad}
\date{August 2026}

\begin{document}
\maketitle

\begin{abstract}
\noindent
Parkinson's disease is classified into Tremor Dominant (TD), Indeterminate and
Postural Instability / Gait Difficulty (PIGD) motor subtypes by a fixed ratio of
MDS-UPDRS items. That rule describes a patient's present state and says nothing
about their future. This project builds the probabilistic extension: given what
is known about a patient today, the probability that they occupy each subtype at
one, two, three and five years.

The work proceeds in three parts. A labelled longitudinal panel of 1{,}080
sporadic PD patients was built from the PPMI dataset, with the three-class
subtype recomputed from raw MDS-UPDRS items and validated at 100.00\% against
PPMI's own two-class derivation and 99.87\% against an archived three-class one.
The 3{,}361 observed transitions in that panel were then characterised. Finally
five progression models were fitted and validated on held-out patients.

Three findings shaped everything after them. \textbf{Indeterminate is a crossing,
not a state} --- only 19.2\% of patients remain in it. \textbf{35.2\% of observed
subtype changes reverse at the very next visit}, yet the movement is
significantly asymmetric ($p = 1.8 \times 10^{-8}$), so the data is a genuine
drift seen through a noisy measurement. And \textbf{a Markov chain in the observed
label is provably misspecified} ($\chi^2 = 156.6$, $p = 7.9 \times 10^{-33}$).

No single model wins. A hidden Markov model with misclassification predicts a
patient's next visit best once they have follow-up history, beating the discrete
baseline in 10 of 10 splits. The proposal's hierarchical Bayesian growth model
supplies the interpretable progression rate --- the TD/PIGD ratio declines
$0.093$ per year (95\% CI $0.068$--$0.116$) --- but predicts subtype less well
and is systematically miscalibrated. For a patient seen for the \emph{first}
time, which is the clinical question, both are beaten by a direct model using
baseline measures, at 71.9\% accuracy.

Every objective set out for this phase of the project has been met. The work is
reproducible from a short sequence of commands and is covered by 144 unit tests.
\end{abstract}

\clearpage
\tableofcontents
\clearpage

% =============================================================================
\section{Objective}
% =============================================================================

The TD/PIGD classification (Stebbins et al., 2013) computes two subscores from
sixteen MDS-UPDRS items and takes their ratio:
\[
  \text{ratio} = \frac{\text{mean of 11 tremor items}}{\text{mean of 5 gait and postural items}}
\]
with $\text{ratio} \geq 1.15 \Rightarrow$ TD, $\text{ratio} \leq 0.90 \Rightarrow$
PIGD, and the interval between them Indeterminate.

The rule is widely used and clinically validated, but it is a snapshot. Patients
move between subtypes as the disease progresses, and the classification carries
no statement about how likely that is. The project's objective is to attach
probabilities to it without replacing it --- to keep the existing thresholds and
report, for an individual patient, the chance of remaining where they are or
moving.

\subsection{Scope}

This report covers the first phase of the project: building the dataset,
measuring how subtypes change over time, and fitting and validating models that
turn those measurements into predictions. Section~9 lists what was produced and
where it lives; Section~10 states plainly what has \emph{not} been done.

% =============================================================================
\section{Constructing a labelled longitudinal dataset}
% =============================================================================

\subsection{Sources}

Four files, of which two had to be downloaded separately.

\begin{table}[H]
\centering
\small
\begin{tabular}{@{}lrl@{}}
\toprule
\textbf{Source} & \textbf{Size} & \textbf{Role} \\
\midrule
PPMI Curated Data Cut       & 19{,}450 $\times$ 212 & visit skeleton, covariates \\
MDS-UPDRS Part II (items)   & 37{,}777 $\times$ 22  & 3 of the 16 items \\
MDS-UPDRS Part III (items)  & 39{,}310 $\times$ 65  & 13 items, medication state \\
Archived Apr-2020 cut       & \phantom{0}3{,}445 $\times$ 159 & three-class validation only \\
\bottomrule
\end{tabular}
\end{table}

\begin{keypoint}
\textbf{The curated cut alone cannot support the project}, which was established
before any modelling. It ships \code{pigd} (the five-item PIGD mean) but no
tremor subscore, so the ratio has a denominator and no numerator. Its
\code{td\_pigd} is coded \emph{two}-class (TD / non-TD), collapsing Indeterminate
away entirely. Both the continuous ratio and the middle class had to be
recomputed from raw items.
\end{keypoint}

\subsection{Construction}

\begin{table}[H]
\centering
\begin{tabular}{@{}lrr@{}}
\toprule
\textbf{Step} & \textbf{Rows} & \textbf{Removed} \\
\midrule
Curated cut                                  & 19{,}450 & --- \\
Restrict to sporadic PD                      & \phantom{0}5{,}789 & 13{,}661 \\
Drop unscheduled and withdrawal visits       & \phantom{0}5{,}774 & \phantom{00,0}15 \\
\midrule
MDS-UPDRS Part III, all rows                 & 39{,}310 & --- \\
Keep OFF-medication or untreated exams       & 29{,}356 & \phantom{0}9{,}954 \\
Merge with Part II, deduplicate              & 37{,}886 & \phantom{00,0}11 \\
\midrule
\textbf{Panel after joining}                 & \textbf{5{,}774} & \\
\quad of which carry a subtype label         & \textbf{4{,}441} & (76.9\%) \\
\bottomrule
\end{tabular}
\caption{Every removed row was verified: all 9{,}954 dropped from Part III have
\code{PDSTATE == 'ON'}, and none are OFF or untreated. The OFF exam is used
because levodopa suppresses tremor, so scoring an ON exam would manufacture
transitions that are pharmacological rather than pathological.}
\end{table}

\subsection{Two data traps}

Both would have corrupted results silently rather than raising an error.

\paragraph{``Unable to Rate'' is stored as the number 101.} In all thirteen
Part III items, in the same numeric field as a real 0--4 rating. A single 101
among the five PIGD items lifts their mean to roughly 20, driving the ratio to
zero and relabelling a tremor-dominant visit as PIGD. \code{NP3PSTBL} alone
carries 3{,}206 of them among OFF exams. The pipeline excludes values outside the
valid 0--4 range, so a sentinel PPMI introduces later cannot enter a subscore
either.

\paragraph{The archived validation file uses a different coding.} Its
\code{td\_pigd\_old} defines $2 =$ PIGD and $3 =$ Indeterminate, the reverse of
the natural ordering. Compared without remapping, agreement reads 65.9\% and
looks like a broken classifier. Confirmed against the archived data dictionary,
remapped, and unit-tested.

\subsection{Validation}

\begin{table}[H]
\centering
\begin{tabular}{@{}llr@{}}
\toprule
\textbf{Reference} & \textbf{Classes} & \textbf{Agreement} \\
\midrule
Curated cut \code{td\_pigd}   & 2 (TD / non-TD)        & 4{,}431 / 4{,}431 $=$ \textbf{100.00\%} \\
Archived \code{td\_pigd\_old} & 3 (TD / Indet.\ / PIGD) & 1{,}595 / 1{,}597 $=$ \textbf{99.87\%} \\
\bottomrule
\end{tabular}
\caption{The three-class check covers 390 patients and is the only available test
of the Indeterminate/PIGD boundary, since the curated cut collapses those two.}
\end{table}

\subsection{Result}

\textbf{1{,}080 patients, 5{,}774 visits, 4{,}441 labelled, up to 14 years of
follow-up.} Baseline distribution: TD 71.7\%, Indeterminate 10.1\%, PIGD 18.1\%.

% =============================================================================
\section{What subtype transitions actually look like}
% =============================================================================

3{,}361 consecutive labelled visit pairs from the 934 patients with at least two
labelled visits; median interval 1.00 years.

\begin{table}[H]
\centering
\begin{tabular}{@{}lrrrr@{}}
\toprule
& \textbf{\arr\ TD} & \textbf{\arr\ Indeterminate} & \textbf{\arr\ PIGD} & \textbf{Row $n$} \\
\midrule
TD            & \textbf{0.808} & 0.080 & 0.112 & 2{,}261 \\
Indeterminate & 0.440 & \textbf{0.192} & 0.368 & \phantom{0,}334 \\
PIGD          & 0.184 & 0.114 & \textbf{0.702} & \phantom{0,}766 \\
\bottomrule
\end{tabular}
\caption{Pooled transition matrix. All probabilities carry Wilson score
intervals in \code{outputs/transitions/}.}
\end{table}

\begin{figure}[H]
\centering
\includegraphics[width=0.66\textwidth]{transition_matrix.png}
\caption{The same matrix. Cell counts are printed because row totals differ by
an order of magnitude.}
\end{figure}

\subsection{Four findings}

\paragraph{1. Indeterminate is a crossing, not a state.} Only 19.2\% remain;
44.0\% return to TD and 36.8\% move to PIGD. Expected persistence 1.24 visits
against 5.21 for TD. It behaves exactly as the neutral band of a continuous,
noisily measured quantity would --- which is the proposal's premise, measured
rather than assumed.

\paragraph{2. PIGD is not absorbing.} 29.8\% of PIGD visits are followed by a
non-PIGD assessment, 18.4\% straight back to TD. Much of the literature treats
TD~\arr~PIGD as one-directional; this data does not.

\paragraph{3. A third of observed changes reverse immediately.}

\begin{table}[H]
\centering
\small
\begin{tabular}{@{}llrrr@{}}
\toprule
\textbf{From} & \textbf{To} & \textbf{$n$} & \textbf{Reverted} & \textbf{Rate} \\
\midrule
TD & Indeterminate & 134 & 75 & \textbf{0.560} \\
PIGD & Indeterminate & \phantom{0}54 & 27 & \textbf{0.500} \\
PIGD & TD & 112 & 44 & 0.393 \\
TD & PIGD & 172 & 55 & 0.320 \\
Indeterminate & PIGD & \phantom{0}93 & 17 & 0.183 \\
Indeterminate & TD & 108 & 19 & 0.176 \\
\midrule
\multicolumn{2}{@{}l}{\textbf{Overall}} & \textbf{673} & \textbf{237} & \textbf{0.352} \\
\bottomrule
\end{tabular}
\caption{Share of observed state changes that return to the origin at the very
next visit. Changes \emph{into} Indeterminate reverse most.}
\end{table}

\begin{figure}[H]
\centering
\includegraphics[width=0.80\textwidth]{reversion.png}
\caption{The same rates, ordered. Bar colour marks the state being left. A high
rate means the recorded label is a noisy reading rather than a durable change.}
\end{figure}

\begin{keypoint}
\textbf{But the movement is not symmetric, and that distinction carries the
project.} If the changes were pure noise, forward and reverse counts would
match. TD~\arr~PIGD occurs 253 times against 141 in reverse
($p = 1.8\times10^{-8}$); Indeterminate~\arr~PIGD 123 against 87 ($p = 0.016$).
Of 932 cross-state moves, 80.5\% cancels and 19.5\% is net drift toward PIGD.

The correct reading is \emph{a genuine directional drift observed through a noisy
measurement} --- which is the case \emph{for} a latent-variable model, not
against it.
\end{keypoint}

\paragraph{4. Transition probabilities scale with the gap between visits.}
TD~\arr~PIGD is 0.100 at intervals under 1.5 years and 0.289 beyond four. A
single pooled matrix therefore averages one-year and multi-year transitions in
whatever proportion attendance happened to produce; it describes the study design
as much as the disease.

\begin{figure}[H]
\centering
\includegraphics[width=0.92\textwidth]{transitions_by_interval.png}
\caption{Selected transitions against the elapsed interval, with Wilson 95\%
intervals. Washed-out bars marked $\dagger$ rest on fewer than 30 observations
and should be read as intervals rather than point estimates.}
\end{figure}

\subsection{Where patients actually end up}

The matrix above describes single steps. Following the same patients through
several visits shows something a one-step matrix cannot.

\begin{figure}[H]
\centering
\includegraphics[width=\textwidth]{pathway_flow.png}
\caption{The 274 patients assessed at baseline, year 2 and year 4. Ribbon
colour marks the state being left.}
\end{figure}

\begin{keypoint}
\textbf{Indeterminate looks stable and is not.} Its size barely moves --- 26
patients at baseline, 29 at year 2, 25 at year 4 --- but of the 26 who start
there, \textbf{only 6 are still there two years later}. Eleven have moved to
PIGD and nine back to TD, while 23 different patients have arrived to replace
them.

The pool stays the same size with almost entirely different members. A
cross-sectional count would show a flat 10\% band and suggest a stable clinical
category; the flow shows a doorway.
\end{keypoint}

\begin{figure}[H]
\centering
\includegraphics[width=0.92\textwidth]{prevalence_over_time.png}
\caption{Subtype prevalence across follow-up, with the number of patients
actually assessed each year beneath. TD falls from 71.7\% to 39.1\% and PIGD
rises from 18.1\% to 48.9\% --- but by year 8 fewer than 12\% of the baseline
cohort is still contributing, so this mixes genuine progression with selective
dropout. The lower panel is part of the reading, not a footnote.}
\end{figure}

\subsection{Time to PIGD conversion}

\begin{table}[H]
\centering
\begin{tabular}{@{}lrr@{}}
\toprule
& \textbf{Naive (1 visit)} & \textbf{Sustained (2 visits)} \\
\midrule
At risk (not PIGD at baseline) & 885 & 885 \\
Converted                      & 263 \,(29.7\%) & 127 \,(14.4\%) \\
KM cumulative incidence, year 5 & \textbf{35.3\%} & \textbf{15.3\%} \\
\bottomrule
\end{tabular}
\caption{Requiring a conversion to hold for two consecutive visits eliminates
51.7\% of them. The definition of ``converted'' moves the answer more than most
modelling choices downstream of it.}
\end{table}

\begin{figure}[H]
\centering
\includegraphics[width=0.78\textwidth]{time_to_pigd_conversion.png}
\caption{Kaplan--Meier cumulative incidence under both definitions, truncated at
year 12 where the risk set falls below 25. The at-risk counts are printed
beneath the axis.}
\end{figure}

\subsection{A bias with a known direction}

23.1\% of cohort visits carry no label. Roughly half are exams that were
performed but could not be fully scored --- and in \textbf{1{,}243 of those
1{,}253 exams the missing item is \code{NP3PSTBL}}, item 3.12, postural
stability. The pull test is often unsafe in a patient with poor balance.

Those exams come from materially more advanced patients: median Hoehn \& Yahr
3.0 against 2.0 (Mann--Whitney $p = 1.4\times10^{-17}$). Postural stability is
one of the five \emph{PIGD} items, so the visits lost are disproportionately
those of patients most likely to be classified PIGD.

\begin{keypoint}
Every progression estimate in this report is therefore a \textbf{lower bound}.
Relaxing the all-items rule is not a remedy: allowing one missing item recovers
44 visits and changes no existing label, because the incomplete exams are
missing many items rather than one.
\end{keypoint}

% =============================================================================
\section{Probabilistic models of progression}
% =============================================================================

Five models were fitted and validated on held-out patients. Splitting is always
at the level of the \emph{patient}, stratified by baseline subtype, so no
individual appears on both sides.

\subsection{A Markov chain in the observed label is misspecified}

A Markov chain assumes the current state summarises the past. That is directly
testable on the 2{,}427 labelled triples.

\begin{table}[H]
\centering
\begin{tabular}{@{}lrrr@{}}
\toprule
\textbf{Currently} & \textbf{Triples} & \textbf{$\chi^2$ (dof 4)} & \textbf{$p$} \\
\midrule
TD            & 1{,}583 & 156.6 & $7.9\times10^{-33}$ \\
Indeterminate & \phantom{0,}235 & \phantom{0}15.3 & $4.2\times10^{-3}$ \\
PIGD          & \phantom{0,}609 & \phantom{0}61.2 & $1.7\times10^{-12}$ \\
\bottomrule
\end{tabular}
\end{table}

For visits currently labelled TD, arriving from PIGD raises the probability of
being PIGD next visit from \textbf{0.086 to 0.393}. Where a patient came from
matters enormously, even conditioning on where they are. That is the signature of
a latent state read through noise: if the true phenotype is PIGD but one exam
scores TD, the next exam is likely to score PIGD again.

\subsection{Model comparison}

\begin{table}[H]
\centering
\begin{tabular}{@{}lrrr@{}}
\toprule
\textbf{Model} & \textbf{Log-lik / transition} & \textbf{Accuracy} & \textbf{Brier} \\
\midrule
HMM (latent state $+$ misclassification) & $\mathbf{-0.6384}$ & \textbf{0.770} & \textbf{0.3458} \\
Discrete Markov chain (per visit)        & $-0.6653$ & 0.768 & 0.3625 \\
Continuous-time Markov chain             & $-0.6752$ & 0.766 & 0.3670 \\
Marginal (null)                          & $-0.8633$ & 0.625 & 0.5185 \\
\bottomrule
\end{tabular}
\caption{276 held-out patients, 1{,}002 transitions.}
\end{table}

\begin{keypoint}
The hidden Markov model beats the discrete baseline in \textbf{10 of 10}
independent splits, by $+0.0284 \pm 0.0094$ log-likelihood per transition
(paired $t$-test $p = 5.4\times10^{-6}$, Wilcoxon $p = 0.002$).

The continuous-time chain --- the expected improvement, given that transition
probabilities scale with the interval --- \textbf{does not help}. It loses at
every interval and by most at long ones. Imposing $P(t) = e^{Qt}$ enforces a
consistency across intervals that a noisily observed process does not satisfy,
so a correct-looking constraint on a misspecified model makes it worse.
\end{keypoint}

\begin{figure}[H]
\centering
\includegraphics[width=\textwidth]{model_comparison.png}
\caption{Every model on every split. Grey lines join the four models within one
split, so a line that never crosses shows a consistent ranking.}
\end{figure}

\subsection{What the hidden Markov model says}

\begin{table}[H]
\centering
\begin{tabular}{@{}lrrr@{}}
\toprule
\textbf{From} & \textbf{\arr\ TD} & \textbf{\arr\ Indeterminate} & \textbf{\arr\ PIGD} \\
\midrule
TD            & $-0.0994$ & 0.0937 & \textbf{0.0057} \\
Indeterminate & 0.1156 & $-0.4713$ & 0.3557 \\
PIGD          & \textbf{0.0000} & 0.0605 & $-0.0605$ \\
\bottomrule
\end{tabular}
\caption{Latent generator $Q$, rates per year.}
\end{table}

\paragraph{Progression is sequential.} The direct TD~\arr~PIGD rate is 0.0057
per year against 0.0937 for TD~\arr~Indeterminate --- \textbf{16 times smaller}
--- and direct PIGD~\arr~TD is estimated at exactly zero. The model was never
told the states are ordered; it recovered that. This is what a continuous
underlying ratio implies: to cross from above 1.15 to below 0.90 a patient must
pass through the middle band.

\paragraph{The label is unreliable in a specific place.} The model estimates how
often an examination records each true phenotype correctly:

\begin{table}[H]
\centering
\begin{tabular}{@{}lrrr@{}}
\toprule
\textbf{True state} & \textbf{recorded TD} & \textbf{recorded Indet.} & \textbf{recorded PIGD} \\
\midrule
TD            & \textbf{0.941} & 0.038 & 0.021 \\
Indeterminate & 0.250 & \textbf{0.500} & 0.250 \\
PIGD          & 0.106 & 0.052 & \textbf{0.842} \\
\bottomrule
\end{tabular}
\caption{Estimated misclassification. The Indeterminate row rests on an imposed
floor of 0.50, so it reads as ``at most half of Indeterminate labels are
correct''. This independently corroborates the finding in Section~3 that changes
into Indeterminate reverse most often (50--56\%), by an entirely different route.}
\end{table}

\paragraph{The phenotype is more stable than its label.} Separating measurement
error from genuine change recovers a much slower underlying process:

\begin{table}[H]
\centering
\begin{tabular}{@{}lrr@{}}
\toprule
\textbf{State} & \textbf{Observed label} & \textbf{Latent phenotype} \\
\midrule
TD            & 3.41 & \textbf{10.06} \\
Indeterminate & 0.44 & \phantom{0}\textbf{2.12} \\
PIGD          & 2.28 & \textbf{16.52} \\
\bottomrule
\end{tabular}
\caption{Mean years before leaving a state. Treating the label as the state
makes every phenotype look transient.}
\end{table}

\begin{figure}[H]
\centering
\includegraphics[width=\textwidth]{progression_curves.png}
\caption{Predicted subtype over time by current phenotype. Because the latent
process runs in continuous time, any horizon can be read off.}
\label{fig:progression}
\end{figure}

% =============================================================================
\section{A hierarchical Bayesian growth model of the ratio}
% =============================================================================

Implemented as specified (proposal pp.\ 13--15):
\begin{align*}
  R^{obs}_{it} &= R_{it} + \epsilon^{obs}_{it}, &
  R_{it} &= (\beta_0 + b_{0i}) + (\beta_1 + b_{1i})\, t + \epsilon^{proc}_{it} \\
  b_{0i} &\sim \mathcal{N}(0, \tau_0^2), & b_{1i} &\sim \mathcal{N}(0, \tau_1^2)
\end{align*}
Inference by Gibbs sampling with conjugate updates throughout; roughly one second
over 573 patients, validated by recovering known parameters from simulated data.

\begin{table}[H]
\centering
\begin{tabular}{@{}lrrr@{}}
\toprule
\textbf{Parameter} & \textbf{Mean} & \textbf{2.5\%} & \textbf{97.5\%} \\
\midrule
$\beta_0$ population baseline ratio   & 2.2920 & 2.1577 & 2.4272 \\
$\beta_1$ population progression, /yr & $-0.0927$ & $-0.1159$ & $-0.0682$ \\
$\sigma$ total residual SD            & 1.1380 & 1.0993 & 1.1778 \\
$\tau_0$ SD of patient baselines      & 1.3408 & 1.2466 & 1.4421 \\
$\tau_1$ SD of patient slopes, /yr    & 0.1033 & 0.0780 & 0.1315 \\
\bottomrule
\end{tabular}
\end{table}

\begin{figure}[H]
\centering
\includegraphics[width=\textwidth]{growth_trajectories.png}
\caption{Population trajectory with 60 patients drawn from the fitted
random-effect distribution, and the posterior for the progression rate.}
\end{figure}

\noindent
\textbf{The ratio declines 0.0927 per year}, $P(\beta_1 < 0) = 1.000$, robust
across five very different prior specifications ($\tau_1$ moves only within
0.094--0.116). From the mean baseline of 2.292 the population average reaches the
TD threshold after a median 12.3 years.

\subsection{Two problems in the specification, found by implementing it}

\paragraph{The two variances are not separately identifiable.} Substituting the
growth model into the observation model,
\[
  R^{obs}_{it} = (\beta_0 + b_{0i}) + (\beta_1 + b_{1i})\,t
    + \underbrace{(\epsilon^{proc}_{it} + \epsilon^{obs}_{it})}_{\text{variance } \sigma^2_{proc} + \sigma^2_{obs}}
\]
Both terms are additive on the same observation, so the likelihood depends only
on their sum --- verified numerically. Only $\sigma^2 = \sigma^2_{proc} +
\sigma^2_{obs}$ is estimated, and this is stated rather than a split reported
that the data cannot support. It matters because $P(TD) = P(R^{true} > 1.15)$
refers to the latent ratio, whose spread excludes observation noise.

\paragraph{A Normal likelihood cannot represent the ratio}, which is
non-negative, right-skewed (skew 1.27) and has an atom at exactly zero (7.0\% of
visits). The model places \textbf{12.2\% of its five-year predictive mass below
zero}.

\subsection{Comparison with the hidden Markov model}

\begin{table}[H]
\centering
\begin{tabular}{@{}lrrr@{}}
\toprule
\textbf{Model} & \textbf{Log-lik / transition} & \textbf{Accuracy} & \textbf{Brier} \\
\midrule
HMM (latent states)          & $\mathbf{-0.6559}$ & \textbf{0.761} & \textbf{0.3573} \\
Hierarchical Bayesian growth & $-0.6814$ & 0.737 & 0.3786 \\
\bottomrule
\end{tabular}
\caption{825 shared held-out transitions. Over six independent splits the HMM is
ahead in 6 of 6, mean $+0.0579$ (sd 0.0188), paired $t$-test $p = 6.6\times10^{-4}$.}
\end{table}

\begin{keypoint}
\textbf{The reason is measurable, not a matter of model capacity.} The growth
model's residual SD is $\sigma = 1.138$; the Indeterminate band is $0.25$ wide.
\textbf{The noise is 4.6 times the width of the band the model must predict
membership of.} Even a perfectly estimated trajectory places the next
observation almost anywhere.
\end{keypoint}

\noindent
The two models are complementary rather than competing. Only the growth model
supplies an interpretable progression rate with a credible interval,
patient-specific slopes, and a route to the $\sigma_i$ stability phenotype in the
proposal's second half.

% =============================================================================
\section{The clinical question: a patient seen for the first time}
% =============================================================================

Every model above earns its accuracy by filtering a patient's \emph{history}
forward. A new patient has none. Measured directly, over five splits, predicting
subtype at horizon $h$ from a first visit only:

\begin{table}[H]
\centering
\begin{tabular}{@{}lr@{}}
\toprule
\textbf{Model (baseline information only)} & \textbf{Held-out log-likelihood} \\
\midrule
Baseline subtype $+$ ratio $+$ covariates & $\mathbf{-0.7550}$ \\
$+$ covariates                            & $-0.7625$ \\
$+$ baseline ratio                        & $-0.7735$ \\
Baseline subtype $+$ horizon              & $-0.7813$ \\
Hidden Markov model (nothing to filter)   & $-0.7869$ \\
Carry forward (``assume no change'')      & $-1.4114$ \\
\bottomrule
\end{tabular}
\end{table}

\begin{keypoint}
\textbf{On the first-visit question the hidden Markov model is the worst of the
five}, not because it is a poor model but because it has nothing to work with. A
direct model of $P(\text{subtype at } h \mid \text{first visit})$ beats it by
$+0.032$, at \textbf{71.9\% accuracy}, and is well calibrated across the full
probability range.

The HMM is not superseded --- its advantage over a lookup table grows with
history: $+0.005$ at one visit, $+0.021$ at two, $+0.047$ at six or more. Which
model to use is decided by how much follow-up exists.
\end{keypoint}

\subsection{Example clinical output}

\begin{table}[H]
\centering
\begin{tabular}{@{}lrrr@{}}
\toprule
\textbf{Horizon} & \textbf{P(TD)} & \textbf{P(Indeterminate)} & \textbf{P(PIGD)} \\
\midrule
now      & 1.000 & 0.000 & 0.000 \\
1 year   & 0.821 & 0.094 & 0.086 \\
2 years  & 0.785 & 0.105 & 0.110 \\
3 years  & 0.749 & 0.113 & 0.138 \\
5 years  & 0.671 & 0.122 & 0.207 \\
\bottomrule
\end{tabular}
\caption{A new TD patient with an otherwise average baseline profile. This is
the output format the proposal specifies.}
\label{tab:clinical-new}
\end{table}

\subsection{Which baseline measures matter}

\begin{table}[H]
\centering
\small
\begin{tabular}{@{}lr|lr@{}}
\toprule
\textbf{Measure} & \textbf{Cost of dropping} & \textbf{Measure} & \textbf{Cost of dropping} \\
\midrule
MDS-UPDRS Part II   & $+0.0153$ & MoCA          & $-0.0005$ \\
DaTscan putamen     & $+0.0050$ & UPSIT         & $-0.0008$ \\
Age                 & $+0.0042$ & Disease duration & $-0.0010$ \\
Hoehn \& Yahr       & $+0.0021$ & SCOPA-AUT     & $-0.0014$ \\
Sex                 & $+0.0002$ & RBDSQ         & $-0.0016$ \\
\bottomrule
\end{tabular}
\caption{Held-out cost of removing each covariate, one at a time.}
\end{table}

\begin{keypoint}
\textbf{The four non-motor markers the proposal's second half is built on ---
RBDSQ, SCOPA-AUT, UPSIT and MoCA --- all sit at or below zero.} The covariate
gain comes from the motor and imaging measures.

This is \emph{not} a refutation of the proposal's hypothesis, which concerns
trajectory \emph{stability} ($\sigma_i$) rather than the trajectory mean, and
leave-one-out understates a predictor that is redundant given another. But it is
worth establishing before that half of the project is built on them.
\end{keypoint}

% =============================================================================
\section{Exploratory analysis and model validation}
% =============================================================================

\subsection{The ratio, examined before any model is fitted}

\begin{itemize}
  \item \textbf{Within-patient SD 1.22} against an Indeterminate band 0.25 wide.
        Band membership was always going to be noise-dominated, and this is
        visible before a model is fitted.
  \item \textbf{ICC 0.536} --- 54\% of ratio variance is between patients.
  \item Raw per-patient OLS slopes have SD \textbf{0.531} against the
        hierarchical $\tau_1$ of 0.10, a factor of five. Shrinkage does real
        work, and a prior on $\tau_1$ must not be centred on the naive spread.
\end{itemize}

\paragraph{One result argues against the modelling choice.} Lag-1 residual
autocorrelation is $-0.219$. Read against zero this looks like mean reversion,
but it cannot be read against zero: residuals from a line fitted to a handful of
points are negatively correlated by construction. Simulating iid noise on the
same visit pattern gives a null of $\mathbf{-0.332}$ $[-0.365, -0.294]$.

The observed value is \emph{more positive} than the null and outside it, so
deviations persist between visits --- the signature of accumulating process
noise. That favours the \textbf{state-space} form over the growth form, and the
state-space form is also where $\sigma_{obs}$ and $\sigma_{proc}$ become
separately identifiable. Two independent reasons to prefer it.

\subsection{How much of a patient-specific estimate is real}

\begin{table}[H]
\centering
\begin{tabular}{@{}lrrrrr@{}}
\toprule
\textbf{Visits} & 2 & 3 & 5 & 8 & 10+ \\
\midrule
Patients                & 191 & 233 & 71 & 40 & 50 \\
Median span (yr)        & 1.16 & 2.09 & 4.17 & 8.42 & 11.67 \\
\textbf{Slope reliability}     & \textbf{0.008} & \textbf{0.024} & 0.112 & 0.384 & 0.569 \\
Baseline reliability    & 0.737 & 0.804 & 0.854 & 0.857 & 0.877 \\
\bottomrule
\end{tabular}
\caption{Reliability is $1 - \mathrm{Var}_{post}(b)/\tau^2$: zero means the
estimate is the population mean with the patient contributing nothing.}
\end{table}

\begin{keypoint}
\textbf{Median slope reliability is 0.050 and no patient exceeds 0.8.} For the
typical patient, 95\% of their ``individual progression rate'' is the cohort
average. The \emph{baseline level}, by contrast, is well identified for
essentially everyone (median 0.805).

So the defensible clinical claim is that a prediction is individualised
\textbf{by position, not by rate}: we know where this patient is, and assume they
move at the cohort's pace. Held-out, that individualisation is worth $+0.20$
log-likelihood, of which 95\% comes from the baseline level.
\end{keypoint}

\subsection{Validation by withholding each patient's future visits}

Withholding each patient's future visits, which is the design the proposal
specifies:

\begin{table}[H]
\centering
\begin{tabular}{@{}lrrrr@{}}
\toprule
\textbf{Train on first} & \textbf{Future visits} & \textbf{Log-lik} & \textbf{Accuracy} & \textbf{95\% coverage} \\
\midrule
1 visit  & 2{,}707 & $-1.0341$ & 0.576 & 0.974 \\
\textbf{2 visits} & 1{,}886 & $\mathbf{-0.8917}$ & 0.647 & \textbf{0.929} \\
3 visits & 1{,}256 & $-0.9398$ & 0.666 & 0.905 \\
\bottomrule
\end{tabular}
\end{table}

\noindent
Interval coverage is close to nominal, so the model's stated uncertainty is
honest. Its \emph{calibration} is not: the growth model under-predicts PIGD in
every bin (0.018 against 0.078 observed, rising to 0.535 against 0.729), mean
absolute gap 0.105. The bias is systematic, one-directional, and unique to this
model among the three.

% =============================================================================
\section{Which model answers which question}
% =============================================================================

\begin{table}[H]
\centering
\small
\begin{tabular}{@{}p{6.6cm}ccc@{}}
\toprule
\textbf{Question} & \textbf{First-visit} & \textbf{HMM} & \textbf{Growth} \\
\midrule
New patient, no follow-up                    & \textbf{best} & worst & middle \\
Patient with several visits                  & --- & \textbf{best} & weaker \\
Progression rate on the clinical scale       & no & no & \textbf{yes} \\
Patient-specific progression rate            & no & no & \textbf{yes} \\
Quantifies label misclassification           & no & \textbf{yes} & no \\
Uses visits with an undefined ratio (16.3\%) & \textbf{yes} & \textbf{yes} & no \\
Well calibrated                              & \textbf{yes} & \textbf{yes} & no \\
Route to the $\sigma_i$ phenotype            & no & no & \textbf{yes} \\
\bottomrule
\end{tabular}
\end{table}

\begin{keypoint}
No single model dominates, and the project should keep all three. The clinical
instrument is the first-visit model; the HMM takes over as follow-up accrues;
the growth model is the scientific description of progression and the foundation
for the proposal's second half.
\end{keypoint}

% =============================================================================
\section{Summary of deliverables}
% =============================================================================

\begin{table}[H]
\centering
\small
\begin{tabular}{@{}p{5.4cm}p{9.0cm}@{}}
\toprule
\textbf{What was produced} & \textbf{Where it lives} \\
\midrule
Labelled longitudinal dataset ---
one row per patient-visit with the TD/PIGD ratio and three-class subtype
  & \code{outputs/visit\_panel.csv}, validated at 100.00\% and 99.87\% against
    two independent PPMI derivations \\
\addlinespace
Description of the ratio's behaviour ---
distribution, within- and between-patient variability, effect of irregular
follow-up
  & \code{outputs/exploration/} \\
\addlinespace
Empirical transition analysis ---
matrices with confidence intervals, interval stratification, reversion,
time-to-conversion, pathways
  & \code{outputs/transitions/} (9 tables, 6 figures) \\
\addlinespace
Five fitted progression models with held-out comparison
  & \code{outputs/model/}, \code{outputs/hbayes/}, \code{outputs/newpatient/} \\
\addlinespace
Per-patient latent phenotype and progression rate, each with a credible
interval and a reliability figure
  & \code{outputs/validation/latent\_trajectories.csv} and
    \code{patient\_parameters.csv} \\
\addlinespace
Validation --- held-out patients, withheld future visits, calibration and
interval coverage
  & \code{outputs/validation/}, \code{outputs/model/} \\
\addlinespace
Clinical output tables giving the probability of each subtype at 1, 2, 3 and 5
years
  & Table~\ref{tab:clinical-new} and Figure~\ref{fig:progression} \\
\bottomrule
\end{tabular}
\caption{Everything set out for this phase of the project has been produced.
Each stage also writes a plain-text report recording every number quoted here.}
\end{table}

% =============================================================================
\section{Limitations}
% =============================================================================

\begin{enumerate}
  \item \textbf{Label availability is not random.} 23.1\% of visits are
        unlabelled, concentrated in patients who cannot complete the pull test
        and who are more advanced. All progression estimates are lower bounds.

  \item \textbf{Per-patient progression rates are largely not recoverable.}
        Median slope reliability 0.050; no patient above 0.8. Individualisation
        is by baseline level, not rate.

  \item \textbf{The $\sigma_i$ analysis is not feasible on the full cohort.} A
        patient-specific variance needs more information than a slope, and after
        fitting an intercept and slope the median patient has one residual degree
        of freedom. Only 182 patients have six or more visits.

  \item \textbf{The growth model is systematically miscalibrated}, under-
        predicting PIGD in every bin.

  \item \textbf{Time-homogeneity is assumed and violated.} The TD~\arr~PIGD rate
        is 0.045/yr before year 3 and 0.085/yr after.

  \item \textbf{16.3\% of labelled visits have no finite ratio}, so any
        ratio-based model silently loses them.

  \item \textbf{Sporadic PD only}, and validation is a random split of one
        cohort. Nothing has been tested externally.

  \item \textbf{The emission floor is an assumption.} The HMM's Indeterminate row
        rests on it; unconstrained the MLE is degenerate.
\end{enumerate}

% =============================================================================
\section{Recommendations}
% =============================================================================

\begin{enumerate}
  \item \textbf{Move to a transformed observation scale.} Three symptoms ---
        $\sigma$ exceeding the band width, negative predictions, right skew ---
        all follow from a symmetric Normal on a non-negative quantity. Thresholds
        carry over exactly, since $R > 1.15 \iff \log R > \log 1.15$. Care is
        needed because the ratio is exactly zero at some visits.

  \item \textbf{Adopt the state-space form.} Two independent arguments now point
        to it: residual autocorrelation exceeds its iid null, and it is the form
        in which the two variances become identifiable.

  \item \textbf{Model the two subscores rather than their ratio.} This recovers
        the 16.3\% of visits with no finite ratio and removes the zero-denominator
        problem at source.

  \item \textbf{Settle the $\sigma_i$ question on the 182-patient subcohort}
        before building the second half of the project on it.

  \item \textbf{Report reliability alongside every per-patient estimate.}
\end{enumerate}

% =============================================================================
\section{Reproducibility}
% =============================================================================

\begin{verbatim}
    python scripts/00_data_audit.py          # verify source files
    python scripts/01_build_panel.py         # labelled longitudinal panel
    python scripts/02_transition_analysis.py # observed transition analysis
    python scripts/03_progression_model.py   # Markov and hidden Markov models
    python scripts/04_hierarchical_bayes.py  # hierarchical Bayesian growth model
    python scripts/05_new_patient_model.py   # first-visit prediction
    python scripts/06_ratio_exploration.py   # exploratory analysis of the ratio
    python scripts/07_longitudinal_cv.py     # latent posteriors, longitudinal CV
    python -m pytest tests/ -q               # 144 tests
\end{verbatim}

\noindent
Roughly 3{,}000 lines of Python across fourteen modules, with every item list,
threshold, cohort rule and model parameter in \code{config/config.yaml} and
validated on load. Nothing is hardcoded in the source. Each stage writes its
tables, figures and a plain-text report recording every number quoted here.

\vspace{1em}
\noindent\textbf{Reference.} Stebbins GT, Goetz CG, Burn DJ, Jankovic J, Khoo TK,
Tilley BC. How to identify tremor dominant and postural instability/gait
difficulty groups with the Movement Disorder Society Unified Parkinson's Disease
Rating Scale. \emph{Movement Disorders} 2013;28(5):668--70.

\end{document}
